%!TEX root = iopjournal.tex
%
% Background
%
% Inhalt (Vorschlag):
%     - CD-Taxonomie und Definitionen: Real Concept Drift vs. Covariate Shift,
%       Sudden / Incremental
%     - performance-basierte Detektoren: KSWIN, ADWIN, DDM, EDDM
%     - Adaptionsstrategien: passiv, aktiv, kombiniert; Stability-Plasticity-Dilemma
%     - der TEP als Anlage: Reaktionen, fuenf Hauptkomponenten, Regelkonzept nach
%       Ricker, Betriebskosten als Zielgroesse, Stoerungskatalog IDV(1)-IDV(28)
%     - Rahmen des Benchmarks: der TEP steht stellvertretend fuer die
%       Betriebspunktoptimierung ueber ein Kostensurrogat; Sensoralterung und
%       begrenzte Kalibrierintervalle sind der Anwendungsbezug, auf den die
%       Ergebnisse generalisieren sollen
%
% Vgl. Dissertation: Kapitel 2 und sec:TEP
%

\section{Background} \label{sec:background}


